수능이 끝난 기룡이는 졸업 전에 추억을 쌓기 위해 반 친구들을 불러 모아 홍삼 게임을 하기로 했다.

홍삼 게임은 다음과 같이 진행된다.

\begin{enumerate}
\item 첫 술래가 ``아\textasciitilde{}싸 너 너!''를 외치며 서로 다른 두 사람을 지목한다.
\item 지목된 두 사람은 새 술래가 되어 각자 ``아\textasciitilde{}싸 너 너!''를 외치며 한 사람씩 지목한다.
\item 새로 지목된 사람이 다시 술래가 되어 이를 반복한다.
\item 한 술래가 다른 친구를 지목하지 못해도 게임은 끝나지 않으며, 남은 한 술래만 계속해서 이어간다.
\item 모든 술래가 다른 친구를 지목하지 못하면 게임은 끝나며, 지목하지 못했던 두 술래는 벌칙을 받는다.
\end{enumerate}

그런데 사람들은 저마다 친하다고 생각하는 사람만 지목한다. 
그래서 친구마다 지목할 수 있는 사람이 정해져 있고, 
심지어 지목할 친구가 아무도 없는 친구도 있다.
이대로라면 한 번도 지목받지 못하고 소외되는 친구가 생길 수도 있다.

친구들 사이의 관계를 모두 알고 있는 기룡이는 아무도 소외되지 않기를 바랐다. 
그래서 자신은 게임에서 빠지는 대신, 매 판 처음 두 사람만 지목하기로 했다.
다만 술래들이 실제로 누구를 지목할지는 기룡이도 알 수 없다. 
따라서 모든 친구가 지목받을 가능성이 있도록 두 사람을 고르기로 했다.

기룡이는 매판 신선함을 느끼도록 처음 지목하는 두 사람의 조합을 매번 다르게 고르려고 한다.
이때 순서만 다른 두 조합은 같은 조합이다.

최대한 게임을 많이 하고 싶은 기룡이를 위해, 모든 친구가 지목될 수 있는 판을 최대 몇 번 할 수 있는지 구해 보자.